\section{Discussion and conclusions}
\label{sec:discussion}
\subsection{What the revised simulations establish}
The revised implementation separates whole-hotspot emission, column resistance, chemical storage and the evidence available to an operator. Consistent tile geometry now gives the configured footprint, damage changes the corresponding source cells, and the timeline and spatial source share the same area weighting. Exact initial and final times, event boundaries and frozen-state plume diagnostics remove the previous age offset. These changes materially alter the deployment interpretation of the earlier plots.

High attenuation does not establish high recovery. Physical resistance persists when chemical storage changes little, and nominal capacity may remain unused because accessibility and equilibrium partitioning limit uptake. The numerical studies quantify discretisation effects; they do not make the assumed material parameters measured properties. The corrected barrier reference uses the same discrete stationary operator, avoiding a small spurious chemistry increment from comparing different boundary discretisations.

Material scale remains direct arithmetic: at 4~kg\,m$^{-2}$, every hectare requires 40~t of core. Whether a small treated area could reduce a large share of exposure depends on measured source heterogeneity and receptor connectivity, neither of which the default uniform hotspot supplies.

\subsection{Interpreting bumps and non-monotone plots}
The revised plots preserve discontinuities at prescribed source changes, physical damage and accepted replacement. Before and after states at the same event time are connected vertically. These are consequences of explicit scenario interventions. A source decrease can also lower the attenuation ratio temporarily because stored contaminant and porewater respond on a finite time scale while the bare-reference denominator changes immediately. Such a response need not indicate negative capture or a failed integrator.

The actual source increase in scenario C illustrates the denominator effect in the opposite direction: at year 2, Pb attenuation rises from 82.72\% to 85.08\% while absolute residual emission also rises. The bare-reference flux increases faster than the residual flux at that instant. The subsequent sharper attenuation decline near years 2.2--2.25 coincides with the Pb capacity ceiling being reached, without a service event. In scenario D, displacement at year 1.5 lowers Pb attenuation by 9.55 percentage points. The lost-tile survey arrives later, and replacement on day 750 raises it by 10.89 points while lowering average active-media saturation. A separate local-damage event at year 2.2 lowers attenuation by 3.17 points. Its subsequent punctured survey class does not satisfy the strict integrity-upper-bound replacement threshold, so that damage does not produce a second service. These timings reflect the specified events, survey cadence and policy rules.

Coastal endpoint maxima vary with the duration and tidal phase of the plume window. A 24-hour window is not an integer number of the prescribed M2 cycles. Its final spatial maximum is a phase-dependent statistic, not a monotonic steady-state convergence measure. The same-phase and numerical-resolution studies separate these effects. Maps are displayed with their actual mat age and matched with/without concentration scales; all time curves include the full reported range.

Some sharp responses arise from constitutive assumptions that remain unvalidated. In particular, exactly capacity-locked cells do not desorb in the implemented law, whereas a just-unsaturated cell can respond to flushing. The independent clean-water comparison exposes this discontinuity. It is retained explicitly for this demonstrator, and should be revisited before using the model to predict remobilisation or service life.

\subsection{Limits of the accounting and monitoring models}
The full-footprint reactive-column ledger closes for the numerical column problem and records active and retrieved inventories separately. Whole-hotspot emission includes bypass and changing effective area. These two diagnostics do not constitute a closed, fully coupled sediment--mat--water inventory when bypass or damage changes: the column inventory is not scaled by the bypassed fraction, and no finite sediment reservoir is resolved. Attenuation against a hypothetical bare reference is therefore not synonymous with contaminant recovered in keratin.

The columns use a clean overlying-water boundary over the long timeline; episodic plume windows start clean, hold the source fixed, and do not feed concentrations back into the earlier column history. The two-dimensional model prescribes currents and mixing rather than solving hydrodynamic momentum. These are explicit reduced-model assumptions. A field assessment would need transport connectivity, heterogeneous seepage, realistic boundaries and appropriately coupled budgets.

The observation clock now spans the experiment, campaign draws are independent of the decision partition, analytical fractions are matched, and quality-control flags affect the controller. Zero burial and missing Cu no longer create positive evidence of a failure. Observations from retired media and exposure windows crossing replacement are excluded from current-media estimates. Nevertheless, the instrument layout remains sparse, uncertainty intervals are heuristic rather than coverage-calibrated probabilities, and recommendations for additional sampling do not yet alter the sample programme. A zero-service outcome cannot by itself prove that a mat was healthy or that monitoring saved a lifecycle cost.

The coupled treatment of censoring is also incomplete. Although the data contract preserves non-detect and above-range bounds and the standalone estimator can propagate them, the current scalar QC has no applicable interval check for the generated censored records. It marks them not evaluated, and the passed-only estimation gate excludes them. Scenario A therefore lacks usable Hg chamber evidence throughout its one-year run even though two bounded Hg chamber results have arrived. This conservative information loss helps explain repeated sampling requests; it does not establish absent Hg flux or validated use of non-detect evidence. In scenarios B and C, later quantified chemistry is available, but wide saturation intervals evaluated at their midpoint remain below the replacement trigger. The fully saturated physical Pb state in scenario C is not directly visible to the policy.

\subsection{Material evidence and priorities for experiments}
The supplied papers support keratin as a candidate sorbent and provide numerical laboratory observations. They do not identify a seawater-qualified finished wool/feather core. The reviewed high mercury uptake concerns chemically modified human hair; the high multi-metal removal concerns a keratin--graphene-oxide composite. Batch removal percentage, uptake at one initial concentration and fitted maximum capacity are different quantities. This report preserves those distinctions and retains the demonstration's operating parameters as assumptions.

A useful experimental sequence should test the finished core in representative seawater and porewater, with dry loading, pH, salinity, sulphide, organic matter and metal fractions measured. Batch uptake should be complemented by flow-through breakthrough, desorption, mixed-ion competition and ageing. Empty geotextile and inert-core controls are needed to separate resistance from chemical capture, and an alternative sorbent can test whether keratin adds a practical advantage.

Ecological assessment must include fibre and particle release, degradation products, effects on benthic exchange and potential methylmercury production. The model contains no microbial Hg transformation or verified ecological benefit. Any later pilot should address a characterised authorised source and measurable flux or exposure endpoints, with retrieval and disposal designed alongside installation.

\subsection{Conclusions}
The revised repository provides reproducible synthetic scenarios, consistent spatial and temporal reporting, a documented numerical audit, and a traceable distinction between published material evidence and assumed operating settings. The plot anomalies included both correct responses to interventions or tides and implementation defects; the identified integration and reporting defects have been corrected and the affected results regenerated. The remaining limitations concern material calibration, constitutive choices, coupling, sparse evidence, censoring-aware QC and deployment feasibility. The defensible use is hypothesis testing and experimental design, with field efficacy and reliable servicing advantage still to be established.
